\section{Finite summands and residue supports}\label{sec:preliminaries}

For $C\subseteq\N_0$, put $C(Y)=|C\cap[0,Y]|$.
All logarithms are natural.

We use $e(x)=\exp(2\pi i x)$ and $e_m(x)=e(x/m)$.
An expectation over a finite set means its uniform probability average
unless another probability law is specified. Multiplicative characters,
including the principal character, are extended by zero at zero.
For functions on $\F_p$ our Fourier convention is
\[
 \widehat f(a)=\frac1p\sum_{x\in\F_p}f(x)e_p(-ax),
 \qquad
 f(x)=\sum_{a\in\F_p}\widehat f(a)e_p(ax).
\]
Norms of functions on a finite field use probability measure. A mass
function $\mu$ of a probability measure instead uses the counting norm
$\|\mu\|_2^2=\sum_x|\mu(x)|^2$ when explicitly so stated below.
This distinction is useful in the collision estimate.

\subsection{A sieve for finitely many shifts}

We use the classical additive large sieve in the following form.
If complex coefficients $c_n$ are supported on an interval of $M$
consecutive integers and $Q\ge1$, then
\begin{equation}\label{eq:prelim-large-sieve}
 \sum_{q\le Q}\ \sum_{\substack{h\bmod q\\(h,q)=1}}
 \left|\sum_n c_n e_q(hn)\right|^2
 \ll (M+Q^2)\sum_n|c_n|^2.
\end{equation}
Indeed, the reduced fractions with denominators at most $Q$ are
$Q^{-2}$-separated modulo one, so this is the separated-frequency
inequality of Montgomery and Vaughan~\cite[Theorem~1]{MontgomeryVaughan1973};
see also Green and Harper~\cite[Proposition~3.1]{GreenHarper}.
The estimate applies to arbitrary complex coefficients. Restricting its
outer sum to prime or squarefree moduli is permitted by nonnegativity.

\begin{lemma}[Fixed shifts]\label{lem:shifted-sieve}
Let $b_1,\ldots,b_j$ be distinct nonnegative integers, where $j\ge1$ is fixed.
If $C\subset\N_0$ and each $c+b_i$ is prime for all sufficiently large
$c\in C$, then
\[
 C(Y)\ll_{b_1,\ldots,b_j} \frac{Y}{(\log Y)^j}.
\]
The implicit constant depends only on the shifts; the threshold for $Y$
may also depend on the finite exceptional range.
\end{lemma}

\begin{proof}
Take $Q=\sqrt Y$ and remove from $C\cap[0,Y]$ all $c\le Q+C_0$,
where $C_0$ is a fixed constant covering the exceptional range.
The remaining set $U$ avoids $j$ distinct classes modulo every prime
$p_0<p\le Q$, for a fixed sufficiently large $p_0$ depending on the
shifts. We may assume $U\ne\varnothing$ and $p_0>j$.

At such a prime define a mean-zero function equal to $1$ on the
$p-j$ allowed classes and to $-(p-j)/j$ on the forbidden classes.
Its probability squared norm is $(p-j)/j$. For squarefree $u$ with
prime factors in this range, the tensor product of these functions
has mean $1$ under the projection of the uniform measure on $U$.
Its additive Fourier expansion contains only primitive modes, since
every local factor has mean zero. Parseval and Cauchy--Schwarz therefore
give
\begin{equation}\label{eq:prelim-primitive-lower}
 \sum_{\substack{h\bmod u\\(h,u)=1}}
 \left|\E_{n\in U}e_u(hn)\right|^2
 \ge \prod_{p\mid u}\frac{j}{p-j}.
\end{equation}
The empty product at $u=1$ is included. Applying
\eqref{eq:prelim-large-sieve} with $c_n=\ind_U(n)/|U|$ bounds the sum
of the left side over $u\le Q$ by $O(Y/|U|)$.

To bound the sum of the right side below, first allow every squarefree
product of primes $p_0<p\le Q^{1/(10j)}$. Its total weight is
\[
 \prod_{p_0<p\le Q^{1/(10j)}}\left(1+\frac{j}{p-j}\right)
 \gg_{b_1,\ldots,b_j} (\log Q)^j
\]
by Mertens' estimates. Under the probability law obtained by normalizing
these weights, a prime is included with probability $j/p$, so
\[
 \E\log u
 =j\sum_{p_0<p\le Q^{1/(10j)}}\frac{\log p}{p}
 =\left(\frac1{10}+o(1)\right)\log Q.
\]
Markov's inequality retains a fixed positive proportion of the weight
when $u\le Q$. Thus $|U|\ll_{b_1,\ldots,b_j}Y/(\log Y)^j$.
Restoring $O(\sqrt Y)$ removed points proves the claim.
\end{proof}

\begin{lemma}[Finite summands]\label{lem:finite-factor}
If $A,B\subset\N_0$ have at least two elements each and $A+B$ agrees
with $\mathcal P$ outside a finite set, then both $A$ and $B$ are infinite.
\end{lemma}

\begin{proof}
Suppose $A$ is finite. Two distinct elements of $A$ give two shifts
of $B$ to which Lemma~\ref{lem:shifted-sieve} applies. Hence
$B(Y)\ll Y/(\log Y)^2$. On the other hand, every sufficiently large
prime at most $Y$ is represented using elements of $A\cap[0,Y]$ and
$B\cap[0,Y]$, so
\[
 \frac{Y}{\log Y}\ll A(Y)B(Y)
 \le |A|B(Y)\ll \frac{Y}{(\log Y)^2},
\]
a contradiction. The other case is symmetric.
\end{proof}

\subsection{The two-infinite-set theorem}

It remains to prove the following technical theorem.

\begin{theorem}[Two infinite summands]\label{thm:main}
There do not exist two infinite sets $A,B\subseteq\N_0$ such that
$(A+B)\mathbin{\triangle}\mathcal P$ is finite.
\end{theorem}

Together with Lemma~\ref{lem:finite-factor}, this proves
Theorem~\ref{cor:ostmann}: a counterexample to the latter would have
both summands infinite and would therefore contradict
Theorem~\ref{thm:main}.

For the remainder of the proof, suppose to the contrary that infinite
sets $A,B\subseteq\N_0$ satisfy this eventual equality.  Fix a threshold
$N$ so that both of the following statements hold:
\begin{equation}\label{eq:eventual-equality}
 \mathcal P\cap(N,\infty)\subseteq A+B,
 \qquad (A+B)\cap(N,\infty)\subseteq\mathcal P.
\end{equation}
The first is prime coverage; the second excludes composite sums.
Choose an integer $N_*>|N|+10$ and put $D=-B$.
Constants in $O(\cdot)$ and $\ll$ may depend on these fixed data;
additional dependencies will be indicated where the order of parameters
matters.

\subsection{Complementary residue supports}

For every prime $p$, the residue sets
\[
 \{a\bmod p:a\in A,\ a>p+N_*\},\qquad
 \{d\bmod p:d\in D,\ -d>p+N_*\}
\]
are disjoint. Indeed, a common residue would make $a-d$ a positive
multiple of $p$ larger than $p$ and $N$, although $a-d\in A+B$ must
be prime. Both sets are nonempty by infinitude. Choose a partition
$S_p,S_p^c$ of $\F_p$ containing the respective sets, and write
$\sigma_p=|S_p|/p$. Thus $0<\sigma_p<1$.
Residues in neither tail support may be assigned to either part.

\subsection{Sizes of the summands}

The following bounds are due to Elsholtz
\cite[Theorem~1.9]{Elsholtz2006}.  We include the large-sieve and
collision argument in the present nonnegative-integer convention.

\begin{lemma}[Square-root bounds]\label{lem:size}
Under the assumed decomposition with both summands infinite,
\begin{equation}\label{eq:src1}
 \frac{\sqrt Y}{(\log Y)^3}\ll A(Y),B(Y)
 \ll \sqrt Y(\log Y)^2.
\end{equation}
\end{lemma}

\begin{proof}
Coverage and the prime number theorem give
$A(Y)B(Y)\gg Y/\log Y$. Infinitude and
Lemma~\ref{lem:shifted-sieve} give, for every fixed positive integer $j$,
\[
 A(Y),B(Y)\ll_j\frac{Y}{(\log Y)^j},
 \qquad A(Y),B(Y)\gg_j(\log Y)^{j-1}.
\]
The second assertion follows from the first and the product lower bound.

Suppose now that $m=A(x)\ge3\sqrt x$, with $x$ sufficiently large.
The set $U=A\cap(\sqrt x+N_*,x]$ has size comparable to $m$.
For $\sqrt x/2<p\le\sqrt x$ it avoids all the
$\nu_p=|B\bmod p|$ classes in $-B\bmod p$. Here all elements of $B$
are allowed: if $p\mid a+b$, then $a+b>p,N$, contrary to primality.
Parseval and Cauchy--Schwarz on the allowed classes give the
nonzero-frequency energy lower bound
\[
 \sum_{h\in\F_p^\times}|\E_{a\in U}e_p(ha)|^2
 \ge \frac{\nu_p}{p-\nu_p}\ge\frac{\nu_p}{p}.
\]
The large sieve then implies
$\sum_{\sqrt x/2<p\le\sqrt x}\nu_p/p\ll x/m$.
Consequently Cauchy--Schwarz, followed by the prime number theorem,
gives
\begin{equation}\label{eq:prelim-H-lower}
 H:=\sum_{\sqrt x/2<p\le\sqrt x}\frac{\log p}{\nu_p}
 \gg \frac{m}{\sqrt x\log x}.
\end{equation}

For any $Y\ge2$ with $B(Y)>0$, consider two independent uniform
elements of $B\cap[0,Y]$. Their collision probability modulo $p$ is
at least $1/\nu_p$. Each unequal pair contributes at most $\log Y$
to the sum of $\log p$ over its prime divisors; equal pairs have
probability $1/B(Y)$ and cost $O(\sqrt x)$. Hence
\begin{equation}\label{eq:prelim-H-upper}
 H\le\log Y+O\left(\frac{\sqrt x}{B(Y)}\right).
\end{equation}
If $m>x^{1/2+\epsilon}$ on an unbounded sequence, for a fixed
$\epsilon>0$, choose $\log Y=x^{\epsilon/2}$ in
\eqref{eq:prelim-H-upper}. Equation~\eqref{eq:prelim-H-lower} forces
$B(Y)\ll\sqrt x$, whereas the preceding polylogarithmic lower bound
contradicts this for sufficiently large fixed $j$. Applying the same
argument to $B$, and using coverage, proves
$A(x),B(x)=x^{1/2+o(1)}$.

We can now take $Y=x^2$ in \eqref{eq:prelim-H-upper}, since
$B(x^2)=x^{1+o(1)}$. Thus $H\ll\log x$, and
\eqref{eq:prelim-H-lower} yields
$m\ll\sqrt x(\log x)^2$. If $m<3\sqrt x$ the same bound is immediate.
The bound for $B$ is symmetric. Coverage once again gives both lower
bounds in \eqref{eq:src1}.
\end{proof}

\subsection{Collision stability}

The congruent-pair count underlying Gallagher's larger sieve
\cite{Gallagher1971} also measures how close a summand is to uniform
on its residue support.  This quantitative use of the larger sieve is
central to Green and Harper's inverse-sieve arguments
\cite[Equation~(2.1) and Lemma~2.4]{GreenHarper}.
We need the following weighted form for the two complementary supports.

Write $U_E$ for uniform probability on a finite nonempty set $E$.
If $\mu$ is a probability measure on integers, write $\mu_p$ for its
projection modulo $p$. The norms in the next statement are counting
norms of probability mass functions.

\begin{lemma}[Collision stability]\label{lem:collision}
Let $b\ge1$ be fixed. Suppose $\mu,\nu$ are probability measures on
\[
 A\cap(\sqrt X+N_*,X],\qquad
 D\cap[-X,-\sqrt X-N_*),
\]
respectively, and that every point mass is at most
$(\log X)^b/\sqrt X$. Put $Q=\sqrt X/(\log X)^{b+1}$. Then
\begin{equation}\label{eq:src2}
\begin{split}
 \sum_{p\le Q}\log p\bigg(
 &\|\mu_p-U_{S_p}\|_2^2+\|\nu_p-U_{S_p^c}\|_2^2\\
 &+\frac{\sigma_p^{-1}+(1-\sigma_p)^{-1}-4}{p}
 \bigg)\ll_b\log\log X.
\end{split}
\end{equation}
All terms on the left are nonnegative.
\end{lemma}

\begin{proof}
The same unequal-pair count used above gives
\[
 \sum_{p\le Q}(\log p)\|\mu_p\|_2^2
 \le\log X+O\left(Q\max_n\mu(n)\right)
 =\log X+O(1),
\]
and the corresponding bound holds for $\nu$.
The projected supports lie in $S_p,S_p^c$. Therefore
\[
 \|\mu_p-U_{S_p}\|_2^2=\|\mu_p\|_2^2-\frac1{|S_p|},
 \quad
 \|\nu_p-U_{S_p^c}\|_2^2=\|\nu_p\|_2^2-\frac1{|S_p^c|}.
\]
The left side of \eqref{eq:src2} is consequently the sum of the two
collision energies minus $4\sum_{p\le Q}\log p/p$. Mertens' estimate
gives
\[
 4\sum_{p\le Q}\frac{\log p}{p}
 =2\log X-4(b+1)\log\log X+O(1),
\]
which proves the assertion. Nonnegativity follows also from
$\sigma^{-1}+(1-\sigma)^{-1}\ge4$ for $0<\sigma<1$.
\end{proof}

One useful form of the lemma does not require pointwise bounded test
functions. If $f_p:\F_p\to\C$ has probability squared norm at most one,
then
\begin{equation}\label{eq:prelim-test-stability}
 \left|\E_\mu f_p-\E_{S_p}f_p\right|^2
 \le p\|\mu_p-U_{S_p}\|_2^2,
\end{equation}
by Cauchy--Schwarz. Thus these squared errors have total
$\log p/p$-weighted sum $O_b(\log\log X)$; the same holds on the other
side. In particular this applies to tests bounded by one.
By Lemma~\ref{lem:size}, uniform measures on
\[
 A\cap[X^{9/10},X],\qquad D\cap[-X,-X^{9/10}]
\]
meet the hypotheses for a fixed $b$: their sizes are
$\gg\sqrt X/(\log X)^3$, while the discarded initial pieces have
size $o(\sqrt X/(\log X)^3)$.
